Because $AC$ is a diameter, $AB\perp BC$. Put $B=(0,0),C=(t,0),A=(0,3\sqrt6)$, with $t>0$. The circle has equation $x^2+y^2-tx-3\sqrt6\,y=0$. The area ratio gives $|y_D|=\sqrt6$. The point $D$ lies on the arc $AC$ opposite $B$, so it is on the positive side of $BC$ and $y_D=\sqrt6$. As $CD=3$, $x_D=t\pm\sqrt3$. Substituting into the circle equation gives $\pm t\sqrt3-9=0$, forcing the plus sign and $t=3\sqrt3$. Thus $AC=\sqrt{t^2+54}=9$, and the radius is $\boxed{\dfrac{9}{2}}$.
\begin{asy}
texpreamble("\usepackage{amsmath,amssymb}");
size(240);
defaultpen(fontsize(9pt));
pair B=(0,0),C=(3*sqrt(3),0),A=(0,3*sqrt(6)),D=(4*sqrt(3),sqrt(6)),O=(A+C)/2;draw(circle(O,4.5));draw(A--B--C--D--A);draw(A--C,dashed);draw(B--D,blue);label("A",A,NW);label("B",B,SW);label("C",C,SE);label("D",D,E);label("O",O,N);
\end{asy}
